%%%%%%%%%%%%%%%%%%%%%%%%%%%%%%%%%%%%%%%%%%%%%%%%%%%%%%%%%%%%%%%%%%%%%%%%%%%%%%%%
%2345678901234567890123456789012345678901234567890123456789012345678901234567890
%        1         2         3         4         5         6         7         8

\documentclass[letterpaper, 10 pt, conference]{ieeeconf} 

\IEEEoverridecommandlockouts % This command is only needed if you want to use the \thanks command
\overrideIEEEmargins         % Needed to meet printer requirements.

\pdfminorversion=4

% --- Packages ---
\usepackage[T1]{fontenc}
\usepackage{url}
\usepackage{cite}
\usepackage{graphicx} 
\usepackage{amsmath} 
\usepackage{amssymb}  
\usepackage{tabularx}
\usepackage{float}
\usepackage{subcaption}
\usepackage{booktabs}
\usepackage{array} 
\usepackage{threeparttable}
\usepackage{epsfig} 
\usepackage{mathptmx} 
\usepackage{times} 

% --- Global Table & Float Spacing Adjustments ---
\setlength{\floatsep}{8pt plus 2pt minus 3pt}       
\setlength{\textfloatsep}{10pt plus 2pt minus 4pt}  
\setlength{\intextsep}{8pt plus 2pt minus 3pt}      
\captionsetup{skip=4pt}       

% --- Universal Spacing Reductions for Display Math (\[ ... \]) ---
\AtBeginDocument{
  \setlength{\abovedisplayskip}{2pt plus 1pt minus 1pt}
  \setlength{\belowdisplayskip}{2pt plus 1pt minus 1pt}
  \setlength{\abovedisplayshortskip}{0pt plus 1pt}
  \setlength{\belowdisplayshortskip}{0pt plus 1pt}
}

% --- Global Hierarchical List Adjustments (Bullet-less Native IEEE Fix) ---
% This safely removes the bullet symbols without requiring external packages that crash IEEE templates.
\renewcommand{\labelitemi}{}
\renewcommand{\labelitemii}{}
\renewcommand{\labelitemiii}{}
\renewcommand{\labelitemiv}{}

\title{\LARGE \bf Optimal Control of a Variable-Mass Thrust-Vectoring Lander*}

\author{Nathan Ge$^{1}$% <-this % stops a space
\thanks{*This work was not supported by any organization}% <-this % stops a space
\thanks{$^{1}$Nathan Ge is a student in the Department of Mechanical and Aerospace Engineering,
        University of California, Los Angeles, 405 Hilgard Avenue, Los Angeles, CA 90095
        {\tt\small nzge@g.ucla.edu}} }

\begin{document}

\maketitle
\thispagestyle{empty}
\pagestyle{empty}

%%%%%%%%%%%%%%%%%%%%%%%%%%%%%%%%%%%%%%%%%%%%%%%%%%%%%%%%%%%%%%%%%%%%%%%%%%%%%%%%
\begin{abstract}
This project studies the optimal guidance and control of a planar variable-mass rocket/lander with translational and rotational dynamics, bounded controls, nonlinear coupling, free-final-time optimal control, endpoint manifold constraints, and trajectory tracking. The project combines Pontryagin’s Maximum Principle (PMP), controllability analysis, Linear Quadratic Regulation (LQR), minimum-time optimal control, nonlinear indirect optimal control, and numerical two-point boundary value problem (TPBVP) solvers. The project is inspired by powered-descent guidance and reusable rocket landing systems studied in modern aerospace optimal control literature. All theoretical derivations must be based on PMP. Numerical implementations may use indirect shooting, multiple shooting, continuation or homotopy methods, and indirect/direct hybrid approaches, provided that the optimality conditions are derived from PMP.
\end{abstract}
%%%%%%%%%%%%%%%%%%%%%%%%%%%%%%%%%%%%%%%%%%%%%%%%%%%%%%%%%%%%%%%%%%%%%%%%%%%%%%%%

github repository: \href{link}{https://github.com/nzge/optimal-control-lander}

\section{Problem Formulation}

The state space of the lander can be represented as: 

\[
x = \begin{bmatrix} p_x & p_z & v_x & v_z & \theta & \omega & m \end{bmatrix}^\top, \quad u = \begin{bmatrix} \delta T & \tau \end{bmatrix}^\top
\]

\begin{table}[H]
    \centering
    \label{tab:parameters}
    \renewcommand{\arraystretch}{1.2} 
    \begin{tabular}{@{} l p{0.65\linewidth} @{}}
        \toprule
        Symbol & Description \\
        \midrule
        \multicolumn{2}{@{}l}{\textit{State Variables ($x$)}} \\
        $p_x, p_z$      & Horizontal and vertical position \\
        $v_x, v_z$      & Horizontal and vertical velocity \\
        $\theta$        & Vehicle pitch angle \\
        $\omega$        & Angular velocity (pitch rate) \\
        $m$             & Instantaneous vehicle mass \\
        \midrule
        \multicolumn{2}{@{}l}{\textit{Control Variables ($u$)}} \\
        $T$             & Commanded thrust magnitude \\
        $\tau$          & Commanded control torque \\
        \bottomrule
    \end{tabular}
\end{table}

The nonlinear equations of motion governing the lander, $\dot{x} = f(x, u)$, are defined by its kinematics, gravity, thrust resolution, and mass depletion rate.

\[
    \dot{x} =
    \begin{bmatrix}
        v_x \\
        v_z \\
        \frac{T}{m} \sin(\theta) \\
        \frac{T}{m} \cos(\theta) - g \\
        \omega \\
        \frac{\tau}{I} \\
        -\frac{T}{I_{sp} g_0}
    \end{bmatrix}
\]

\begin{table}[H]
    \centering
    \caption{System Constraints}
    \label{tab:system_constraints}
    \renewcommand{\arraystretch}{1.2}
    \begin{tabular}{@{} l p{0.65\linewidth} @{}}
        \toprule
        Mathematical Constraint & Description \\
        \midrule
        \multicolumn{2}{@{}l}{\textit{Control Constraints}} \\
        $T \in [T_{\text{min}}, T_{\text{max}}]$ & Actuator bounds on thrust magnitude \\
        $\tau \in [\tau_{\text{min}}, \tau_{\text{max}}]$ & Actuator bounds on control torque \\
        \midrule
        \multicolumn{2}{@{}l}{\textit{Numerically Enforced State Constraints}} \\
        $p_z(t) \geq 0$ & Hard deck limit preventing subsurface flight \\
        $m(t) \geq m_{\text{dry}}$ & Minimum mass limit (dry mass) \\
        \bottomrule
    \end{tabular}
\end{table}

\begin{table}[H]
    \centering
    \caption{System Parameters}
    \label{tab:system_parameters}
    \renewcommand{\arraystretch}{1.2}
    \begin{tabular}{@{} l c p{0.5\linewidth} @{}}
        \toprule
        Symbol & Value & Description \\
        \midrule
        $g$ & $9.81 \text{ m/s}^2$ & Gravitational acceleration \\
        $I$ & $10 \text{ kg}\cdot\text{m}^2$ & Vehicle moment of inertia \\
        $m_0$ & $20 \text{ kg}$ & Initial vehicle mass \\
        $\alpha$ & $0.005 \text{ kg/(N}\cdot\text{s)}$ & Propellant consumption rate \\
        $T_{\text{min}}$ & $0 \text{ N}$ & Minimum thrust limit \\
        $T_{\text{max}}$ & $300 \text{ N}$ & Maximum thrust limit \\
        $\tau_{\text{min}}$ & $-30 \text{ N}\cdot\text{m}$ & Minimum torque limit \\
        $\tau_{\text{max}}$ & $30 \text{ N}\cdot\text{m}$ & Maximum torque limit \\
        \bottomrule
    \end{tabular}
\end{table}

%%%%%%%%%%%%%%%%%%%%%%%%%%%%%%%%%%%%%%%%%%%%%%%%%%%%%%%%%%%%%%%%%%%%%%%%%%%%
\section{Numerical Methodology}

\subsection{Linear Systems Solver: The Sweep Method (Parts I \& II)}
$scipy.integrate.solve\_ivp$ integrates the Riccati equation backward in time to generate the feedback gain matrix $K(t)$. An interpolator was then used on this gain during the forward-pass simulation of the state trajectory.

\subsection{Nonlinear TPBVP Solvers (Parts III \& IV)}
The Riccati substitution no longer applies once the optimal control problem is nonlinear, so the coupled state/costate system arising from PMP is solved directly. Part~III uses \emph{indirect single shooting} on the free-final-time mission: the $2n$ canonical equations are integrated forward and a Newton iteration adjusts the unknown initial costate, the phase-switch time, and the final time until the manifold, transversality, and $H(t_f)=0$ residuals vanish. Part~IV uses \emph{indirect collocation} (\texttt{scipy.integrate.solve\_bvp}) on the fixed-time problem: the $14$ state/costate equations are discretized on a mesh and the optimal control is substituted in closed form from the stationarity conditions. Collocation is preferred for Part~IV because it accepts a full state/costate \emph{trajectory} as its initial guess, which is taken from the corresponding linearized LQR solution; this warm start lies close to the nonlinear extremal and lets the solver converge in a few mesh-refinement iterations. In both parts the actuator bounds are imposed by saturating the PMP-optimal control, consistent with Pontryagin's pointwise minimization of $H$ over the admissible control set.

\subsection{Constraint Enforcement Strategies}

\begin{itemize}
    \item Control Constraints ($T, \tau$): For the linear LQR phases I applied post-calculation saturation, clipping the commanded $\delta T$ and $\tau$ to $[T_{\min},T_{\max}]$ and $[\tau_{\min},\tau_{\max}]$. For the nonlinear PMP phases the control follows in closed form from the stationarity condition $\partial H/\partial u = 0$ and is then clipped to the same actuator bounds inside the solver right-hand side; this is exactly Pontryagin's pointwise minimization of $H$, since for the control-affine Hamiltonian the constrained minimizer is the unconstrained stationary control projected onto the admissible box.
    \item State Constraints ($p_z \ge 0$, $m \ge m_{dry}$): Adjoining state constraints directly to the PMP Hamiltonian creates severe analytical difficulties (jump discontinuities in the costates). The numerical workaround is to either:
    \begin{itemize}
        \item "Soft Constraints" (adding massive penalties to your $Q$ matrix if the rocket goes underground)
        \item "Event Tracking" (terminating the ODE solver upon impact)
        \item Direct constraint arrays in an NLP solver.
    \end{itemize}
\end{itemize}

%%%%%%%%%%%%%%%%%%%%%%%%%%%%%%%%%%%%%%%%%%%%%%%%%%%%%%%%%%%%%%%%%%%%%%%%%%%%%%%%

\newpage
\section{Part I — Fixed-Time LQR Around Different Nominal Trajectories}

\subsection{Phase Formulation}

Linearization: The linearized model plant is formulated as $\dot{x}=Ax+Bu$ where $A = \frac{\partial f}{\partial x}$, $B = \frac{\partial f}{\partial u}$, and $f(x,u) = \dot{x}$.

\begin{itemize}
    \item Hover
\[
\setlength{\arraycolsep}{2.5pt}
\renewcommand{\arraystretch}{0.9}
A_{\text{hover}} = \begin{bmatrix} 
0 & 0 & 1 & 0 & 0 & 0 & 0 \\ 
0 & 0 & 0 & 1 & 0 & 0 & 0 \\ 
0 & 0 & 0 & 0 & g & 0 & 0 \\ 
0 & 0 & 0 & 0 & 0 & 0 & -\frac{g}{m_0} \\ 
0 & 0 & 0 & 0 & 0 & 1 & 0 \\ 
0 & 0 & 0 & 0 & 0 & 0 & 0 \\ 
0 & 0 & 0 & 0 & 0 & 0 & 0 
\end{bmatrix}, \quad
B_{\text{hover}} = \begin{bmatrix} 
0 & 0 \\ 
0 & 0 \\ 
0 & 0 \\ 
\frac{1}{m_0} & 0 \\ 
0 & 0 \\ 
0 & I^{-1} \\ 
-\alpha & 0 
\end{bmatrix}
\]

\item Descent
\[
\setlength{\arraycolsep}{2.5pt}
\renewcommand{\arraystretch}{0.9}
A(t) = \begin{bmatrix} 
0 & 0 & 1 & 0 & 0 & 0 & 0 \\ 
0 & 0 & 0 & 1 & 0 & 0 & 0 \\ 
0 & 0 & 0 & 0 & g & 0 & 0 \\ 
0 & 0 & 0 & 0 & 0 & 0 & -\frac{g}{m^*(t)} \\ 
0 & 0 & 0 & 0 & 0 & 1 & 0 \\ 
0 & 0 & 0 & 0 & 0 & 0 & 0 \\ 
0 & 0 & 0 & 0 & 0 & 0 & 0 
\end{bmatrix}, \quad
\begin{aligned}
&B(t) = \begin{bmatrix} 
0 & 0 \\ 
0 & 0 \\ 
0 & 0 \\ 
\frac{1}{m^*(t)} & 0 \\ 
0 & 0 \\ 
0 & 1/I \\ 
-\alpha & 0 
\end{bmatrix}
\end{aligned}
\]

$m^*(t) = m_0 e^{-\alpha g t} $
\end{itemize}

\subsection{Analytical Derivations (PMP)}

\begin{itemize}
    \item Hamiltonian construction
    \[H(x, u, p) = L(x, u) + p^\top f(x, u)\] 
    \vspace{2pt}\centerline{where $L(x,u)=x^\top Qx + u^\top Ru$}\vspace{2pt}
    \[H(x, u, p) = x^\top Q x + u^\top R u + p^\top (Ax + Bu)\]
    
    \item Transversality
    \[p(t_f) = \frac{\partial \phi}{\partial x}\bigg|_{t=t_f}\]
    \[p(t_f) = 2 Q_f x(t_f) \]
    
    \item Adjoint Equations (costate dynamics)
    \[\dot{p} = -\frac{\partial H}{\partial x}\]
    Using matrix calculus (and the symmetry of $Q$):
    \[\dot{p} = -\big( 2 Q x + A^\top p \big) = -2 Q x - A^\top p \]
    
    \item Stationarity Conditions
    \[\frac{\partial H}{\partial u} = 0 \implies 2 R u + B^\top p = 0 \implies u^* = -\frac{1}{2} R^{-1} B^\top p \]
    
    \item Riccati Differential equation
    Assumption (Riccati Substitution): 
    \[p(t) = P(t)x(t)\] 
    Where $P(t)$ is an unknown, time-varying symmetric matrix. Differentiate the assumption using the product rule:
    \[\dot{p} = \dot{P}x + P\dot{x}\]
    Substitute $\dot{p}$ and $\dot{x}$ (the dynamics).
    \[-2Qx - A^\top P x = \dot{P}x + P(Ax + Bu^*)\]
    Substitute the optimal control $u^*$. Remembering that $p = Px$, it becomes $u^* = -\frac{1}{2} R^{-1} B^\top P x$.
    \[-2Qx - A^\top P x = \dot{P}x + P\left(Ax - \frac{1}{2} B R^{-1} B^\top P x\right)\]
    \[-2Qx - A^\top P x = \dot{P}x + PAx - \frac{1}{2} P B R^{-1} B^\top P x\]
    Group all terms multiplying $x(t)$. Move everything to one side of the equation:
    \[\left[ \dot{P} + P A + A^\top P - \frac{1}{2} P B R^{-1} B^\top P + 2Q \right] x = 0\]
    Because this equation must hold true for any arbitrary state trajectory $x(t)$, the bracketed matrix expression must identically equal zero. Setting the bracket to zero yields the Riccati Differential Equation (RDE): 
    \[\boxed{\dot{P} = -P A - A^\top P + \frac{1}{2} P B R^{-1} B^\top P - 2Q}\]
    
    \item Optimal State-feedback law
    \[P(t_f) = 2 Q_f, \quad u^*(t) = -K(t) x(t), \quad K(t) = \frac{1}{2} R^{-1} B^\top P(t)\]
\end{itemize}

\subsection{Numerical Implementation and LQR Tuning}

\begin{itemize}
    \item $Q$, $R$ and $Q_f$ matrix comparisons. We used heuristic methods to create a selection of $Q$, $R$, and $Q_f$ that accurately reflect underlying control intuition. Fig.~\ref{fig:placeholder_sweep} sweeps these presets on the hover LTI plant: \texttt{pos\_heavy} and \texttt{att\_heavy} tighten position and attitude transients at the cost of higher $\int u^\top R u\,dt$; \texttt{cheap\_th} and \texttt{exp\_trq} isolate thrust- versus torque-dominated effort; \texttt{balanced} is the baseline used in later parts.
\end{itemize}

\begin{table}[H]
  \centering
  \caption{Part~I.5.5 cost presets (diagonal $Q$, $R$, $Q_f$).}
  \label{tab:part1-cost-compact}
  \small
  \setlength{\tabcolsep}{2.5pt} 
  \resizebox{\linewidth}{!}{% 
    \begin{tabular}{@{}l ccc@{}} 
      \toprule
      Preset & $Q=\mathrm{diag}(\cdot)$ & $R$ & $Q_f=\mathrm{diag}(\cdot)$ \\
      \midrule
      balanced        & $(1,1,0.5,0.5,10,5,0.01)$       & $(0.1,1)$  & $(10,10,1,1,50,20,0.1)$ \\
      pos\_heavy      & $(20,20,2,2,5,2,0.01)$          & $(0.1,1)$  & $(50,50,5,5,20,10,0.1)$ \\
      att\_heavy      & $(1,1,0.5,0.5,50,25,0.01)$      & $(0.1,2)$  & $(10,10,1,1,100,50,0.1)$ \\
      cheap\_th       & $(1,1,0.5,0.5,10,5,0.01)$       & $(0.01,1)$ & $(10,10,1,1,50,20,0.1)$ \\
      exp\_trq        & $(1,1,0.5,0.5,10,5,0.01)$       & $(0.1,10)$ & $(10,10,1,1,50,20,0.1)$ \\
      scaled\_id      & $(10^{-2},10^{-2},0.04,0.04,25,4,1)$ & $(1,1)$ & same as $Q$ \\
      \bottomrule
    \end{tabular}% 
  } 
\end{table}

\begin{figure}[H]
    \centering
    \includegraphics[width=\linewidth]{figures/p1_cost_preset_sweep.png}
    \caption{Cost sweep}
    \label{fig:placeholder_sweep}
\end{figure}

\subsection{Analysis}

Controllability Analysis

\begin{itemize}
    \item Algebraic Controllability Matrix (Frozen-Time):
    Used to check the structural control authority of the LTI system (or a "frozen" snapshot of an LTV system). For a system with $n = 7$ states:
    \[C = \begin{bmatrix} B & AB & A^2B & A^3B & A^4B & A^5B & A^6B \end{bmatrix}\]
    
    \item Finite-Horizon Controllability Gramian:
    \[W_c(0, t_f) = \int_{0}^{t_f} \Phi(0, \tau) B(\tau) B(\tau)^\top \Phi(0, \tau)^\top d\tau\]
    Where $\Phi(t, \tau)$ is the state transition matrix of $A(t)$.

    \item Numerical Computation (Differential Equation):
    To practically compute $W_c$ numerically and avoid explicitly calculating the integral and state transition matrix, we can solve its differential equation form:
    \[\dot{W}_c(t) = A(t)W_c(t) + W_c(t)A(t)^\top + B(t)B(t)^\top\]
    
    \item Minimum Control Energy:
    The minimum control energy $E$ required to reach a specific state deviation $x_f$ is given by:
    \[E = x_f^\top W_c^{-1} x_f\]
\end{itemize}


\begin{figure}[H]
    \centering
    \includegraphics[width=\linewidth]{figures/p1_gramian_remaining_horizon.png}
    \caption{Controllability Gramian and horizontal contrallability proxy}
\end{figure}

Controllable energy in the weakest meaningful mode shrinks as 
$t_f - t$ shrinks, especially for horizontal motion (as shown by the proxy). Both trims are algebraically full rank ($7/7$), but the descent LTV plant at $t=15$~s shows larger $B$-column entries for $v_z$ and weaker indirect authority on $p_x$ (horizontal acceleration enters only through attitude coupling), which is the physical origin of the ``weakly controllable'' label.


\begin{itemize}
    \item Controllability matrices
        \begin{itemize}
            \item State order: $x=[p_x,p_z,v_x,v_z,\theta,\omega,m]^\top$
            \item controls $u=[\delta T,\tau]^\top$
            \item $\mathcal{C}=[B,AB,\ldots,A^{6}B]\in\mathbb{R}^{7\times 14}$.
        \end{itemize}
        
        \begin{itemize}
            \item Trim Snapshot $t=0.0\,\mathrm{s}$
            
            Descent trim: $m=20.00\,\mathrm{kg}$, $v_x^\ast=40.0\,\mathrm{m/s}$, $v_z^\ast=-8.0\,\mathrm{m/s}$.
            
            \begin{table}[H]
                \centering
                \caption{Controllability rank and singular values}
                \small
                \setlength{\tabcolsep}{4pt}
                \begin{tabular}{@{}lcccc@{}}
                \toprule
                Model & $\mathrm{rank}(\mathcal{C})$ & $\sigma_1$ & $\sigma_6$ & $\sigma_7$ \\
                \midrule
                Hover LTI & 7/7 & 0.98 & 0.0489 & 0 \\
                Descent LTV & 7/7 & 0.98 & 0.0489 & 0 \\
                \bottomrule
                \end{tabular}
            \end{table}
                
            Hover LTI at $t=0.0\,\mathrm{s}$.
                
            \begin{small}
                \setlength{\arraycolsep}{1.8pt} 
                \renewcommand{\arraystretch}{0.95}
                \noindent\resizebox{\linewidth}{!}{$\displaystyle 
                \begin{array}{lc@{\hspace{2pt}}cc@{\hspace{2pt}}cc@{\hspace{2pt}}cc@{\hspace{2pt}}cc@{\hspace{2pt}}cc@{\hspace{2pt}}cc@{\hspace{2pt}}c}
                 & \multicolumn{2}{c}{B} & \multicolumn{2}{c}{AB} & \multicolumn{2}{c}{A^2B} & \multicolumn{2}{c}{A^3B} & \multicolumn{2}{c}{A^4B} & \multicolumn{2}{c}{A^5B} & \multicolumn{2}{c}{A^6B} \\
                \hline
                p_x & 0 & 0 & 0 & 0 & 0 & 0 & 0 & 0.98 & 0 & 0 & 0 & 0 & 0 & 0 \\
                p_z & 0 & 0 & 0.05 & 0 & 0.00245 & 0 & 0 & 0 & 0 & 0 & 0 & 0 & 0 & 0 \\
                v_x & 0 & 0 & 0 & 0 & 0 & 0.98 & 0 & 0 & 0 & 0 & 0 & 0 & 0 & 0 \\
                v_z & 0.05 & 0 & 0.00245 & 0 & 0 & 0 & 0 & 0 & 0 & 0 & 0 & 0 & 0 & 0 \\
                \theta & 0 & 0 & 0 & 0.1 & 0 & 0 & 0 & 0 & 0 & 0 & 0 & 0 & 0 & 0 \\
                \omega & 0 & 0.1 & 0 & 0 & 0 & 0 & 0 & 0 & 0 & 0 & 0 & 0 & 0 & 0 \\
                m & -0.005 & 0 & 0 & 0 & 0 & 0 & 0 & 0 & 0 & 0 & 0 & 0 & 0 & 0
                \end{array}$}
            \end{small}
                
            Descent LTV at $t=0.0\,\mathrm{s}$.
                
            \begin{small}
                \setlength{\arraycolsep}{1.8pt}
                \renewcommand{\arraystretch}{0.95}
                \noindent\resizebox{\linewidth}{!}{$\displaystyle 
                \begin{array}{lc@{\hspace{2pt}}cc@{\hspace{2pt}}cc@{\hspace{2pt}}cc@{\hspace{2pt}}cc@{\hspace{2pt}}cc@{\hspace{2pt}}cc@{\hspace{2pt}}c}
                 & \multicolumn{2}{c}{B} & \multicolumn{2}{c}{AB} & \multicolumn{2}{c}{A^2B} & \multicolumn{2}{c}{A^3B} & \multicolumn{2}{c}{A^4B} & \multicolumn{2}{c}{A^5B} & \multicolumn{2}{c}{A^6B} \\
                \hline
                p_x & 0 & 0 & 0 & 0 & 0 & 0 & 0 & 0.98 & 0 & 0 & 0 & 0 & 0 & 0 \\
                p_z & 0 & 0 & 0.05 & 0 & 0.00245 & 0 & 0 & 0 & 0 & 0 & 0 & 0 & 0 & 0 \\
                v_x & 0 & 0 & 0 & 0 & 0 & 0.98 & 0 & 0 & 0 & 0 & 0 & 0 & 0 & 0 \\
                v_z & 0.05 & 0 & 0.00245 & 0 & 0 & 0 & 0 & 0 & 0 & 0 & 0 & 0 & 0 & 0 \\
                \theta & 0 & 0 & 0 & 0.1 & 0 & 0 & 0 & 0 & 0 & 0 & 0 & 0 & 0 & 0 \\
                \omega & 0 & 0.1 & 0 & 0 & 0 & 0 & 0 & 0 & 0 & 0 & 0 & 0 & 0 & 0 \\
                m & -0.005 & 0 & 0 & 0 & 0 & 0 & 0 & 0 & 0 & 0 & 0 & 0 & 0 & 0
                \end{array}$}
            \end{small}
        
            \item Trim Snapshot $t=15.0\,\mathrm{s}$
                
            Descent trim: $m=9.59\,\mathrm{kg}$, $v_x^\ast=25.0\,\mathrm{m/s}$, $v_z^\ast=-8.0\,\mathrm{m/s}$.
                
            \begin{table}[H]
                \centering
                \caption{Controllability rank and singular values}
                \small
                \setlength{\tabcolsep}{4pt}
                \begin{tabular}{@{}lcccc@{}}
                \toprule
                Model & $\mathrm{rank}(\mathcal{C})$ & $\sigma_1$ & $\sigma_6$ & $\sigma_7$ \\
                \midrule
                Hover LTI & 7/7 & 0.98 & 0.0489 & 0 \\
                Descent LTV & 7/7 & 0.98 & 0.1 & 0 \\
                \bottomrule
                \end{tabular}
            \end{table}
                
            Hover LTI at $t=15.0\,\mathrm{s}$.
                
            \begin{small}
                \setlength{\arraycolsep}{1.8pt}
                \renewcommand{\arraystretch}{0.95}
                \noindent\resizebox{\linewidth}{!}{$\displaystyle 
                \begin{array}{lc@{\hspace{2pt}}cc@{\hspace{2pt}}cc@{\hspace{2pt}}cc@{\hspace{2pt}}cc@{\hspace{2pt}}cc@{\hspace{2pt}}cc@{\hspace{2pt}}c}
                 & \multicolumn{2}{c}{B} & \multicolumn{2}{c}{AB} & \multicolumn{2}{c}{A^2B} & \multicolumn{2}{c}{A^3B} & \multicolumn{2}{c}{A^4B} & \multicolumn{2}{c}{A^5B} & \multicolumn{2}{c}{A^6B} \\
                \hline
                p_x & 0 & 0 & 0 & 0 & 0 & 0 & 0 & 0.98 & 0 & 0 & 0 & 0 & 0 & 0 \\
                p_z & 0 & 0 & 0.05 & 0 & 0.00245 & 0 & 0 & 0 & 0 & 0 & 0 & 0 & 0 & 0 \\
                v_x & 0 & 0 & 0 & 0 & 0 & 0.98 & 0 & 0 & 0 & 0 & 0 & 0 & 0 & 0 \\
                v_z & 0.05 & 0 & 0.00245 & 0 & 0 & 0 & 0 & 0 & 0 & 0 & 0 & 0 & 0 & 0 \\
                \theta & 0 & 0 & 0 & 0.1 & 0 & 0 & 0 & 0 & 0 & 0 & 0 & 0 & 0 & 0 \\
                \omega & 0 & 0.1 & 0 & 0 & 0 & 0 & 0 & 0 & 0 & 0 & 0 & 0 & 0 & 0 \\
                m & -0.005 & 0 & 0 & 0 & 0 & 0 & 0 & 0 & 0 & 0 & 0 & 0 & 0 & 0
                \end{array}$}
            \end{small}
                
            Descent LTV at $t=15.0\,\mathrm{s}$.
                
            \begin{small}
                \setlength{\arraycolsep}{1.8pt}
                \renewcommand{\arraystretch}{0.95}
                \noindent\resizebox{\linewidth}{!}{$\displaystyle 
                \begin{array}{lc@{\hspace{2pt}}cc@{\hspace{2pt}}cc@{\hspace{2pt}}cc@{\hspace{2pt}}cc@{\hspace{2pt}}cc@{\hspace{2pt}}cc@{\hspace{2pt}}c}
                 & \multicolumn{2}{c}{B} & \multicolumn{2}{c}{AB} & \multicolumn{2}{c}{A^2B} & \multicolumn{2}{c}{A^3B} & \multicolumn{2}{c}{A^4B} & \multicolumn{2}{c}{A^5B} & \multicolumn{2}{c}{A^6B} \\
                \hline
                p_x & 0 & 0 & 0 & 0 & 0 & 0 & 0 & 0.98 & 0 & 0 & 0 & 0 & 0 & 0 \\
                p_z & 0 & 0 & 0.104 & 0 & 0.00511 & 0 & 0 & 0 & 0 & 0 & 0 & 0 & 0 & 0 \\
                v_x & 0 & 0 & 0 & 0 & 0 & 0.98 & 0 & 0 & 0 & 0 & 0 & 0 & 0 & 0 \\
                v_z & 0.104 & 0 & 0.00511 & 0 & 0 & 0 & 0 & 0 & 0 & 0 & 0 & 0 & 0 & 0 \\
                \theta & 0 & 0 & 0 & 0.1 & 0 & 0 & 0 & 0 & 0 & 0 & 0 & 0 & 0 & 0 \\
                \omega & 0 & 0.1 & 0 & 0 & 0 & 0 & 0 & 0 & 0 & 0 & 0 & 0 & 0 & 0 \\
                m & -0.005 & 0 & 0 & 0 & 0 & 0 & 0 & 0 & 0 & 0 & 0 & 0 & 0 & 0
                \end{array}$}
            \end{small}
        \end{itemize}

    \item Controllable and weakly controllable linearization comparisons.

        \begin{figure}[H]
            \centering
            \includegraphics[width=\linewidth]{figures/p1_ic_sensitivity.png}
            \caption{Sensitivity to initial condition error}
            \label{fig:ic_sensitivity}
        \end{figure}
        
        At small perturbations: descent can be worse (weak plant + wrong linearization point).
        At large perturbations: hover overtakes — hover LTI may be overly optimistic at large $\alpha$ (linear model stays fixed while state explores larger deviations).
        
        \begin{figure}[H]
            \centering
            \includegraphics[width=\linewidth]{figures/p1_hover_vs_descent_2x2.png}
            \caption{State and control effort across trajectory}
        \end{figure}
        
        \begin{figure}[H]
            \centering
            \includegraphics[width=\linewidth]{figures/p1_weak_direction_transients.png}
            \caption{Transient behavior}
        \end{figure}

        $p_x$ and $\theta$ swings indicate structural difficulty. Even “controllable” hover struggles. Gramian analysis flagged as expensive to correct. The gap is clearest in $p_z$ (mass/thrust schedule, changing trim). Descent worse

        \begin{figure}[H]
            \centering
            \includegraphics[width=\linewidth]{figures/p1_state_control_histories.png}
            \caption{Optimal solution trajectories versus time: states and control histories.}
            \label{fig:p1_state_hist}
        \end{figure}

        \begin{figure}[H]
            \centering
            \includegraphics[width=0.85\linewidth]{figures/p1_mission_plane_quiver.png}
            \caption{Mission-plane trajectory ($p_x$ vs $p_z$) with attitude vectors at sampled positions.}
            \label{fig:p1_mission_plane}
        \end{figure}

        \subsection{Validation}

        \begin{figure}[H]
            \centering
            \includegraphics[width=\linewidth]{figures/p1_hamiltonian.png}
            \caption{Hamiltonian evolution}
        \end{figure}
        
        Because this phase utilizes the Riccati sweep method rather than a shooting solver, optimality is validated by verifying the backward-pass terminal conditions and forward-pass tracking convergence. Table \ref{tab:val_p1} confirms the Riccati matrix $P(t)$ strictly satisfies the terminal boundary condition $P(t_f) = Q_f$. Additionally, the closed-loop tracking error $e(t)$ is bounded within acceptable numerical tolerances at the final time.
        
        \begin{table}[H]
        \centering
        \footnotesize
        \caption{Riccati and Tracking Validation (Phase I/II)}
        \label{tab:val_p1}
        \begin{tabular}{@{}lllc@{}}
        \toprule
        Metric & Target & Computed & Error (Residual) \\ \midrule
        $||P(t_f) - Q_f||_\infty$ & $0$ & $1.22 \times 10^{-14}$ & $1.22 \times 10^{-14}$ \\
        $||e(t_f)||_\infty$ & $0$ & $4.51 \times 10^{-6}$ & $4.51 \times 10^{-6}$ \\
        \bottomrule
        \end{tabular}
        \end{table}
        
\end{itemize}
        

%%%%%%%%%%%%%%%%%%%%%%%%%%%%%%%%%%%%%%%%%%%%%%%%%%%%%%%%%%%%%%%%%%%%%%%%%%%%%%%%
\newpage
\section{PART II - Optimal Trajectory Tracking}

\subsection{Phase formulation}

Objective (free final state, no terminal cost, as specified for Part~II):
\[
    J = \int_{0}^{t_f} \Big[ (x - x_{\text{ref}})^\top Q (x - x_{\text{ref}}) + u^\top R u \Big] dt
\]

\subsection{Scenario and Operating Point}

The reference $x_{\text{ref}}(t)$ is one of the optimal regulation trajectories recorded in Part I (Fig.~\ref{fig:p2_ref_overview}): a damped maneuver that settles to the hover trim. To produce a nontrivial transient, the controller is initialized at a \emph{substantially different} state, $x_0^{\text{trk}} = [-12,\,15,\,-4,\,3,\,-0.25,\,0.4,\,-0.8]^\top$, deviating from the reference start by $\Delta p_x \approx 17$~m and $\Delta\theta \approx 0.33$~rad. Because the hover linearization is altitude-invariant ($p_z$ is decoupled and absent from every right-hand side), the maneuver is executed about a hover point $h = 12$~m above the pad; this leaves $A$, $B$, $x_{\text{ref}}$, and the feedback law unchanged while the absolute altitude $p_z = h + p_z^{\text{dev}}$ is used to enforce the ground constraint $p_z \geq 0$. The numerically enforced state constraints $p_z \geq 0$ and $m \geq m_{\text{dry}}$ are satisfied throughout: the minimum absolute altitude is $11.6$~m (reference) and $10.0$~m (tracker), and the mass margin remains above $12$~kg.

\subsection{Analytical Derivations (PMP)}

\begin{itemize}

    \item The Hamiltonian:
    \[H = (x - x_{\text{ref}})^\top Q (x - x_{\text{ref}}) + u^\top R u + p^\top (Ax + Bu)\]
    
    \item Adjoint Equation (Costate Dynamics):
    \[\dot{p} = -\frac{\partial H}{\partial x}\] 
    \[\dot{p} = -2Q(x - x_{\text{ref}}) - A^\top p \]
    
    \item Stationarity (Optimal Control Law):
    \[\frac{\partial H}{\partial u} = 0 \implies 2Ru + B^\top p = 0 \implies u^* = -\frac{1}{2} R^{-1} B^\top p\]
    
    \item Transversality (Terminal Condition):

        With a free final state and no terminal cost ($\phi \equiv 0$), the transversality condition gives a homogeneous terminal costate:
    \[p(t_f) = \frac{\partial \phi}{\partial x} \bigg|_{t=t_f} = 0 \]
    
    \item Tracking Riccati Equations (Riccati substitution)
    
        The added auxiliary vector $s(t)$ accounts for the tracking offset:
        \[p(t) = P(t)x(t) + s(t)\]
        \[\dot{p} = \dot{P}x + P\dot{x} + \dot{s}\]
    
        Substitute $\dot{p}$ and $\dot{x}$:
        \[-2Q(x - x_{\text{ref}}) - A^\top(Px + s) = \dot{P}x + P(Ax + Bu^*) + \dot{s}\]
    
        Substitute $u$ ($u^* = -\frac{1}{2}R^{-1}B^\top(Px + s)$):
        \[
            \begin{aligned}
            -2Qx &+ 2Qx_{\text{ref}} - A^\top P x - A^\top s ={} \\
            &\dot{P}x + PAx - \tfrac{1}{2}P B R^{-1} B^\top (Px + s) + \dot{s}
            \end{aligned}
        \]
        
        Group the terms by $x(t)$:
        \[
            \begin{split}
            &\Big[ \dot{P} + PA + A^\top P - \tfrac{1}{2}PBR^{-1}B^\top P + 2Q \Big] x(t) \\
            &\quad + \Big[ \dot{s} + A^\top s - \tfrac{1}{2}PBR^{-1}B^\top s - 2Qx_{\text{ref}} \Big] = 0
            \end{split}
        \]
    
        Since the expression must hold true for any arbitrary trajectory $x(t)$, we get the following tracking Riccati pair:
    
        \begin{itemize}
            \item  Matrix Riccati Differential Equation (The Feedback Rulebook):
            \[\dot{P} = -PA - A^\top P + \frac{1}{2}PBR^{-1}B^\top P - 2Q\]
            
            \item Auxiliary Vector Differential Equation (The Feedforward Tracker):
            \[\dot{s} = - \left( A^\top - \frac{1}{2}PBR^{-1}B^\top \right) s + 2Q x_{\text{ref}}\]
        \end{itemize}

        Since $p(t_f) = P(t_f)x(t_f) + s(t_f) = 0$ must hold for an arbitrary terminal state, both pieces are integrated backward from the homogeneous terminal conditions
        \[P(t_f) = 0, \qquad s(t_f) = 0,\]
        reflecting the absence of any terminal penalty in the Part~II objective.
        
    
    \item Closed-Loop Tracking Control Law

        Once you integrate both $\dot{P}$ and $\dot{s}$ backward from $t_f$ to $0$, you can construct your final control law. Substitute $p(t) = P(t)x(t) + s(t)$ back into $u$: 
        \[u^*(t) = -\frac{1}{2} R^{-1} B^\top P(t) x(t) - \frac{1}{2} R^{-1} B^\top s(t)\]
        This is elegantly split into two components:
        \[u^*(t) = -K(t)x(t) + u_{ff}(t)\]
        $K(t)x(t)$ is the feedback component that stabilizes the rocket and kills deviations. $u_{ff}(t)$ is the feedforward component that anticipates the control effort.
    
\end{itemize}

\subsection{Analysis}

\begin{itemize}
    \item \begin{figure}[H]
            \centering
            \includegraphics[width=\linewidth]{figures/p2_tracking_states_4x2.png}
            \caption{Tracking Performance}
            \label{fig:tracking_perf}
        \end{figure}
        
        Tracker rejoins the reference qualitatively. Errors shrink and shapes align by mid-horizon.

        \begin{figure}[H]
            \centering
            \includegraphics[width=\linewidth]{figures/p2_reference_states_controls.png}
            \caption{Recorded Part I optimal trajectory used as the reference $x_{\text{ref}}(t)$, $u_{\text{ref}}(t)$.}
            \label{fig:p2_ref_overview}
        \end{figure}

        Quantitatively, the tracker starts with state error $\lVert e(0)\rVert \approx 19.9$, peaks at $\lVert e\rVert \approx 26.2$ near $t = 1.6$~s as it first accelerates toward the reference, and the cost-weighted error decays to within $10\%$ of its peak by $t \approx 6.9$~s. The residual terminal error $\lVert e(t_f)\rVert \approx 1.38$ is concentrated in the vertical and mass channels ($|e_{p_z}| \approx 1.16$, $|e_{v_z}| \approx 0.67$, $|e_m| \approx 0.32$); the horizontal and attitude errors are driven to $\mathcal{O}(10^{-3})$. This residual is an expected feature of the fixed-horizon, free final-state problem with no terminal cost ($P(t_f)=0$): as $t \to t_f$ the Riccati gain $K(t)$ and feedforward $s(t)$ both decay, so the controller ``relaxes'' near $t_f$ and the weakly penalized vertical/mass states drift slightly.

        The required mission-plane views (Figs.~\ref{fig:p2_ref_plane}--\ref{fig:p2_conv_plane}) plot $p_x$ versus $p_z$ with body-axis attitude arrows. The tracker spirals down from its displaced initial state and converges onto the reference at the hover target, remaining well above the ground line throughout.

        \begin{figure}[H]
            \centering
            \includegraphics[width=0.78\linewidth]{figures/p2_reference_mission_plane.png}
            \caption{Reference mission trajectory in the $p_x$--$p_z$ plane with attitude.}
            \label{fig:p2_ref_plane}
        \end{figure}

        \begin{figure}[H]
            \centering
            \includegraphics[width=0.82\linewidth]{figures/p2_mission_plane_convergence.png}
            \caption{Closed-loop convergence of the tracker onto the reference, with attitude arrows and the ground constraint $p_z = 0$.}
            \label{fig:p2_conv_plane}
        \end{figure}

    \item Transient Response

        \begin{figure}[H]
            \centering
            \includegraphics[width=\linewidth]{figures/p2_error_transients.png}
            \caption{Transient Response}
            \label{fig:transient_resp1}
        \end{figure}

        Scalar tracking error (cost-balanced error) decays monotonically in the weighted norm. Moving target, not zero, so no rebound mid-horizon. Position errors decay slowly and dominate; attitude error stays small throughout. Playing ‘catch-up’.

        \begin{figure}[H]
            \centering
            \includegraphics[width=\linewidth]{figures/p2_control_decomposition.png}
            \caption{Transient Control Decomposition}
            \label{fig:transient_resp2}
        \end{figure}

        The control splits as $u^*(t) = -K(t)\,e(t) + u_{ff}(t)$. The feedback term dominates the transient and drives the large initial error to zero, while the feedforward term $u_{ff} = -R^{-1}B^\top s$ is small and slowly varying because the reference is itself a decaying signal. The torque channel uses nearly the full authority at $t=0$ ($\tau \approx 29.6$~N$\cdot$m against the $\pm 30$~N$\cdot$m bound) to arrest the attitude error, but no channel saturates over the horizon.

    \item Control Effort

        \begin{figure}[H]
            \centering
            \includegraphics[width=\linewidth]{figures/p2_control_effort.png}
            \caption{Control Effort}
            \label{fig:control_effort}
        \end{figure}

        Relative to the smooth reference plan, rejoining from the displaced initial state is expensive: the tracker expends $\int u^\top R u\,dt \approx 1041$ versus $\approx 86$ for the reference, a $12\times$ increase concentrated in the first few seconds. Peak thrust deviation reaches $\delta T \approx 92$~N (within the $+104$~N headroom to $T_{\max}$), confirming that the chosen weights keep the catch-up maneuver feasible without thrust saturation. Once the error is nulled, the tracking control collapses onto the reference control.

    \item Robustness to Deviations from the Nominal Trajectory

        \begin{figure}[H]
            \centering
            \includegraphics[width=\linewidth]{figures/p2_ic_robustness.png}
            \caption{Robustness to initial-condition deviations: terminal error and tracking cost as the initial offset is scaled by $\alpha$.}
            \label{fig:robustness}
        \end{figure}

        Sweeping the initial offset by a factor $\alpha$ confirms the structure of the linear tracking law: the terminal error grows essentially \emph{linearly} in $\alpha$ (slope $\approx 1.18$), consistent with superposition of the closed-loop transition map, while the tracking cost grows \emph{quadratically} (the cost is a quadratic form in the deviation). There is no threshold or instability over the tested range, demonstrating graceful degradation with respect to initial-condition error.

        Robustness to errors in the assumed reference is examined in Fig.~\ref{fig:p2_ref_noise}, where the reference used to build the feedforward is corrupted by zero-mean noise of level $\sigma$ and the resulting controller is scored against the \emph{true} reference. The terminal error is only mildly affected (from $1.38$ at $\sigma = 0$ to $1.50$ at $\sigma = 0.30$) because the feedback term rejects high-frequency reference error, but the tracking cost degrades markedly (from $\approx 3150$ to $\approx 6810$, a $2.2\times$ increase) as the controller wastes effort chasing the corrupted plan.

        \begin{figure}[H]
            \centering
            \includegraphics[width=\linewidth]{figures/p2_reference_noise_robustness.png}
            \caption{Sensitivity to a perturbed reference plan: terminal error and cost (scored against the true reference) versus reference noise level $\sigma$.}
            \label{fig:p2_ref_noise}
        \end{figure}

\end{itemize}

Compared with Part~I regulation, the tracking problem replaces $Q_f x(t_f)$ with $P(t_f)=s(t_f)=0$ and adds the feedforward $u_{ff}(t)$; the feedback structure is unchanged, but the controller must expend far more effort to ``catch up'' from a displaced initial state.

\subsubsection*{Unknown Reference Trajectory}

The present law is non-causal: $s(t)$ comes from a backward integration that requires the entire future $x_{\text{ref}}(\tau)$ on $[t,t_f]$. Without that preview, only $u=-K(t)e(t)$ remains---stable, but lagging and costlier. Fig.~\ref{fig:p2_ref_noise} shows that reference error propagates into $s(t)$ and inflates cost even when feedback preserves terminal accuracy. A causal fix would replan online (e.g.\ receding-horizon MPC) or generate $x_{\text{ref}}$ from a predictive model.

\subsection{Validation}

Because the Part~II objective has no terminal cost and the final state is free, optimality is certified against the PMP necessary conditions rather than terminal convergence. The tracking Riccati pair satisfies the homogeneous free-terminal conditions $P(t_f)=0$ and $s(t_f)=0$, and the stationarity condition $Ru + B^\top \lambda = 0$ (with costate $\lambda = P(t)\,e(t) + s(t)$) holds along the entire closed-loop rollout. Table~\ref{tab:val_p2} reports these residuals at the level of machine precision. The terminal tracking error is \emph{not} required to vanish for this free-final-state problem; its value ($\lVert e(t_f)\rVert_\infty \approx 1.16$, dominated by the vertical/mass relaxation analyzed above) is a result, not a residual.

\begin{table}[H]
\centering
\footnotesize
\caption{Part II Tracking: PMP Optimality Residuals}
\label{tab:val_p2}
\begin{tabular}{@{}lllc@{}}
\toprule
Metric & Target & Computed & Error (Residual) \\ \midrule
$\|P(t_f)\|_\infty$ (Riccati terminal) & $0$ & $1.0 \times 10^{-14}$ & $1.0 \times 10^{-14}$ \\
$\|s(t_f)\|_\infty$ (feedforward terminal) & $0$ & $1.6 \times 10^{-17}$ & $1.6 \times 10^{-17}$ \\
$\max_t \|Ru + B^\top \lambda\|$ (stationarity) & $0$ & $1.1 \times 10^{-14}$ & $1.1 \times 10^{-14}$ \\
\bottomrule
\end{tabular}
\end{table}

The state-equation defect of the RK4 rollout is $\mathcal{O}(10^{-1})$ at the chosen step. Figure~\ref{fig:p2_hamiltonian} shows the Hamiltonian evolution: for the autonomous reference plan $H$ is constant to $\sim 10^{-4}$, whereas for the tracking problem $H$ is intentionally time-varying because the explicit dependence on $x_{\text{ref}}(t)$ makes $\partial H/\partial t \neq 0$ -- the relevant optimality certificate for tracking is the stationarity/costate residual rather than constancy of $H$.

\begin{figure}[H]
    \centering
    \includegraphics[width=\linewidth]{figures/p2_hamiltonian.png}
    \caption{Hamiltonian consistency: constant for the autonomous reference plan and time-varying for the tracking problem.}
    \label{fig:p2_hamiltonian}
\end{figure}

%%%%%%%%%%%%%%%%%%%%%%%%%%%%%%%%%%%%%%%%%%%%%%%%%%%%%%%%%%%%%%%%%%%%%%%%%%%%%%%%
\newpage
\section{PART III - The Mission Problem}

\subsection{Phase A — Minimum-Time Ascent}

\subsubsection*{Formulation}
\[J_A = \int_{0}^{t_1} 1 \, dt = t_1\]

\subsubsection*{Analytical Derivations (PMP)}

Hamiltonian:
\[H_A(x, u, p) = 1 + p^\top(Ax + Bu) = 1 + p^\top Ax + p^\top Bu\]

Adjoint equation:
\[\dot{p} = -\frac{\partial H_A}{\partial x} = -A^\top p\]

Switching Function:
\[\frac{\partial H_A}{\partial u} = B^\top p \implies S(t) = B^\top p(t)\]

\[
u_i^*(t) = \begin{cases} 
U_{\max, i} & \text{if } S_i(t) < 0 \\
U_{\min, i} & \text{if } S_i(t) > 0 
\end{cases}
\]

No singular arcs arise on thrust ($S_{\delta T}$ has one sign change). The torque channel is structurally singular for this vertical ascent ($S_\tau=p_\omega/I\equiv0$ with $\theta=\omega=0$), so $\tau^\star=0$ throughout.

\[H_A(t) \equiv 0 \implies 1 + p(t)^\top \big(Ax(t) + Bu^*(t)\big) = 0 \quad \forall t \in [0, t_1]\]

\subsubsection*{Endpoint and Free-Final-Time Conditions}

The ascent drives the vehicle to a \emph{prescribed} rendezvous state. The first six components of the state are fixed at $t_1$, while the terminal mass is left free, so the corresponding costate vanishes by transversality:
\[
\begin{aligned}
x_i(t_1) &= x_{sky,i}, \quad i \in \{p_x, p_z, v_x, v_z, \theta, \omega\}, \\
p_m(t_1) &= 0 \quad \text{(free terminal mass)}.
\end{aligned}
\]
Because $t_1$ is free, the minimum-time transversality condition closes the system with $H_A(t_1)=0$. The decision variables $\big(p(0)\in\mathbb{R}^7,\ t_1\big)$ are matched against the eight residuals $\{x_{1:6}(t_1)-x_{sky,1:6},\ p_m(t_1),\ H_A(t_1)\}$.

\subsubsection*{Implementation}

Phase~A is solved by indirect single shooting (\texttt{mission.py}). The decision vector is $Z=[p(0);\,t_1]\in\mathbb{R}^8$; each candidate is propagated forward with \texttt{solve\_ivp} (RK45) on the $14$-dimensional state--costate system using the bang--bang law above, and the eight residuals $\{x_{1:6}(t_1)-x_{sky,1:6},\,p_m(t_1),\,H_A(t_1)\}$ are driven below $10^{-5}$ by \texttt{scipy.optimize.least\_squares} with $12$ random multi-starts and $t_1\in[0.3,\,5]$~s. Controls enter the right-hand side as saturated PMP minimisers; $\tau^\star=0$ because $S_\tau$ is singular for the decoupled vertical ascent. The rendezvous is $x_{sky}=(0,\,8,\,0,\,v_{z,sky},\,0,\,0,\,0)^\top$ (hover-deviation coordinates) with $v_{z,sky}\in\{-2,0,+2\}$~m/s for the required terminal-velocity comparison.

\subsubsection*{Ascent Trajectory Results}

Fig.~\ref{fig:p3_phaseA_vz} shows bang--bang thrust with exactly one switch in every case: maximum $\delta T$ accelerates upward until $S_{\delta T}=(B^\top p)_1$ changes sign, then minimum thrust completes the rendezvous. Thrust saturates because time appears linearly in $J_A$ with no control penalty. Ascent time depends strongly on $v_{z,sky}$ ($t_1=2.39$, $2.16$, and $1.98$~s for $-2$, $0$, and $+2$~m/s); the switch occurs late in the arc ($\approx59$--$71\%$ of $t_1$), so deceleration to the target speed is back-loaded. Propellant use is nearly unchanged ($\approx2.08$~kg) because the maneuver is short; mass depletion is too small over $\sim2$~s to materially change $1/m$ and control authority.

\begin{figure}[H]
    \centering
    \includegraphics[width=\linewidth]{figures/p3_phaseA_terminal_velocity.png}
    \caption{Phase~A minimum-time ascent for terminal vertical velocities $v_{z,sky}\in\{-2,0,+2\}$~m/s: altitude $p_z(t)$, vertical velocity $v_z(t)$, bang--bang thrust deviation $\delta T(t)$, and switching function $S_{\delta T}=(B^\top p)_1$. Each case exhibits a single thrust switch; the switch time, total ascent time $t_1$, and propellant use vary with $v_{z,sky}$.}
    \label{fig:p3_phaseA_vz}
\end{figure}

\subsubsection*{Validation}

\begin{figure}[H]
    \centering
    \includegraphics[width=\linewidth]{figures/p3_switching_controls.png}
    \caption{Phase~A switching function $S_{\delta T}(t)$ (top) and full mission control histories $\delta T(t),\ \tau(t)$ (bottom) for M1 and M2; the dotted line marks the Phase~A/Phase~B handoff $t_1$. The single sign change of $S_{\delta T}$ confirms the bang--bang structure with one switch.}
    \label{fig:p3_switching}
\end{figure}

\subsection{Phase B — Return-and-Land Problem}

Phase B begins precisely where Phase A ends, thus the initial condition for the landing phase is $x(0) = x_1$.

\subsubsection*{Formulation \& Core PMP} 
\[J_B = \int_{0}^{t_f} \Big( x^\top Q x + u^\top R u \Big) dt\]

Generic Hamiltonian, and the fundamental $H(t_f) = 0$ condition:
\[H_B(x, u, p) = x^\top Q x + u^\top R u + p^\top (Ax + Bu)\]

The unconstrained control law and adjoint equations remain identical to Phase I:
\[u^*(t) = -\frac{1}{2} R^{-1} B^\top p(t)\]
\[\dot{p}(t) = -2Qx(t) - A^\top p(t)\]

Because $t_f$ is free, PMP imposes the free-time transversality condition. The Hamiltonian evaluated at the terminal time must equal zero: 

\[
\begin{aligned}
H_B(t_f) = 0 \implies & x(t_f)^\top Q x(t_f) + u(t_f)^\top R u(t_f) \\
& + p(t_f)^\top \big( Ax(t_f) + Bu(t_f) \big) = 0
\end{aligned}
\]

\subsubsection*{Transversality} 

\[p(t_f) = \frac{\partial \phi}{\partial x}\bigg|_{t_f} + \sum_{j} \nu_j \nabla h_j(x(t_f))\]

\begin{itemize}
    \item \textit{Manifold 1 (Pinpoint Landing)}

    \begin{itemize}
        \item Description
        \[
        \begin{aligned}
        h_1(x) &= p_z = 0 \\
        h_2(x) &= v_z = 0 \\
        h_3(x) &= \theta = 0
        \end{aligned}
        \]

        \item Gradients
        \[
        \begin{aligned}
        \nabla h_1 &= [0, 1, 0, 0, 0, 0, 0]^\top \\
        \nabla h_2 &= [0, 0, 0, 1, 0, 0, 0]^\top \\
        \nabla h_3 &= [0, 0, 0, 0, 1, 0, 0]^\top
        \end{aligned}
        \]

        \item Transversality
        \[
        \begin{aligned}
        p_{p_x}(t_f) &= \frac{\partial \phi}{\partial p_x} && \text{(Free state)} \\
        p_{p_z}(t_f) &= \frac{\partial \phi}{\partial p_z} + \nu_1 && \text{(Constrained)} \\
        p_{v_x}(t_f) &= \frac{\partial \phi}{\partial v_x} && \text{(Free state)} \\
        p_{v_z}(t_f) &= \frac{\partial \phi}{\partial v_z} + \nu_2 && \text{(Constrained)} \\
        p_{\theta}(t_f) &= \frac{\partial \phi}{\partial \theta} + \nu_3 && \text{(Constrained)} \\
        p_{\omega}(t_f) &= \frac{\partial \phi}{\partial \omega} && \text{(Free state)} \\
        p_{m}(t_f) &= \frac{\partial \phi}{\partial m} && \text{(Free state)}
        \end{aligned}
        \]
        \end{itemize}
 
    \item \textit{Manifold 2 (Surface Landing)}

        \begin{itemize}
            \item Description: landing on the surface of a dome/sphere.
            \[
            \begin{aligned}
            h_1(x) &= (p_x - p_c)^2 + p_z^2 - r^2 = 0 \\
            h_2(x) &= v_z = 0 \\
            h_3(x) &= \theta = 0
            \end{aligned}
            \]

            \item The Gradients: $h_1$ couples $p_x$ and $p_z$, spanning 2 dimensions.
            \[
            \begin{aligned}
            \nabla h_1 &= [2(p_x - p_c), 2p_z, 0, 0, 0, 0, 0]^\top \\
            \nabla h_2 &= [0, 0, 0, 1, 0, 0, 0]^\top \\
            \nabla h_3 &= [0, 0, 0, 0, 1, 0, 0]^\top
            \end{aligned}
            \]

            \item Transversality
            \[
            \begin{aligned}
            p_{p_x}(t_f) &= \frac{\partial \phi}{\partial p_x} + 2\nu_1(p_x(t_f) - p_c) \\
            p_{p_z}(t_f) &= \frac{\partial \phi}{\partial p_z} + 2\nu_1 p_z(t_f)
            \end{aligned}
            \]
            \textit{(The remaining 5 costates are identical to M1)}

            \item Conclusion for M2: We can eliminate the unknown multiplier $\nu_1$ by dividing the first two equations (assuming $\frac{\partial \phi}{\partial x} = 0$). This yields the exact boundary condition passed to the TPBVP solver:
            \[\frac{p_{p_x}(t_f)}{p_{p_z}(t_f)} = \frac{p_x(t_f) - p_c}{p_z(t_f)}\]
        \end{itemize}  
\end{itemize}

\subsubsection*{Numerical Implementation}

The free-final-time landing problem is solved as a true indirect PMP boundary value problem by collocation. The combined state--costate system $(\dot{x}=Ax+Bu^\star,\ \dot{p}=-2Qx-A^\top p)$ with the stationarity law $u^\star=\mathrm{sat}\!\big(-\tfrac12 R^{-1}B^\top p\big)$ is mapped onto $s\in[0,1]$ via $t=t_1+s\,(t_f-t_1)$, and the unknown duration $t_f-t_1$ enters as a free parameter. The split boundary conditions are the handoff $x(t_1)=x_1$, the seven manifold endpoint and transversality residuals (Section~7.3: $\{p_z,v_z,\theta,p_{p_x},p_{v_x},p_\omega,p_m\}$ for M1; $\{h_1,v_z,\theta,p_{v_x},p_\omega,p_m,\text{collinearity}\}$ for M2), and the free-time closure $H_B(t_f)=0$, giving the fifteen conditions required for the fourteen states plus the parameter. The platform is centred at $(p_c,0)=(0,0)$ so that M1 (flat ground) lands at $(0,0)$ and M2 (circular platform) lands at the apex $(0,r)$; both share the same horizontal footprint, isolating the effect of the terminal-manifold geometry. Among the multiple $t_f$ local optima satisfying the necessary conditions, the global minimiser of $J_B$ that respects $p_z\ge0$ throughout is retained.

\begin{figure}[H]
    \centering
    \includegraphics[width=0.49\linewidth]{figures/p3_trajectory_M1.png}
    \hfill
    \includegraphics[width=0.49\linewidth]{figures/p3_trajectory_M2.png}
    \caption{Mission-plane trajectories $p_x$ vs $p_z$ with body-attitude vectors at sampled positions: M1 flat-ground landing (left, touchdown at the origin) and M2 circular-platform landing (right, touchdown at the platform apex). The ascent (Phase~A) reaches $x_{sky}$ (star); the descent (Phase~B) tracks the optimal landing arc.}
    \label{fig:p3_traj}
\end{figure}

\begin{figure}[H]
    \centering
    \includegraphics[width=0.6\linewidth]{figures/p3_trajectory_M1_M2_overlay.png}
    \caption{Overlaid M1 and M2 mission planes. Both scenarios share the launch site, $x_{sky}$, and horizontal footprint; they differ only in the terminal manifold (flat ground at $(0,0)$ vs.\ circular platform apex at $(0,r)$).}
    \label{fig:p3_overlay}
\end{figure}

\subsubsection*{Landing Manifold Comparison}

Both missions share Phase~A, $x_{sky}$, and the horizontal footprint; only the Phase~B terminal manifold differs (Fig.~\ref{fig:p3_overlay}). Fig.~\ref{fig:p3_comparison} and Table~\ref{tab:val_p3} summarise the five comparison axes:

\begin{itemize}
    \item \textit{Landing time.} M2 completes the mission sooner: $t_f=7.32$~s versus $8.10$~s for M1 ($\Delta t_f\approx0.78$~s, Table~\ref{tab:val_p3}). Phase~B alone is also shorter ($5.16$ vs.\ $5.94$~s) because touchdown at the platform apex $(0,r)=(0,5)$~m avoids the long final vertical flare required to reach $p_z=0$ on flat ground.

    \item \textit{Energy consumption.} Despite the shorter clock time, M2 costs more control energy: $\int_{t_1}^{t_f} u^\top R u\,dt \approx 526$ versus $216$ for M1 ($\approx2.4\times$, Table~\ref{tab:val_p3}). The lofted descent arc and stronger horizontal--vertical coordination dominate the energy budget.

    \item \textit{Switching structure.} Phase~A is identical in both cases (one bang--bang thrust switch at $t_1=2.16$~s). Phase~B uses smooth saturated costate feedback in both manifolds---no further bang--bang arcs. Fig.~\ref{fig:p3_switching} shows similar overall control profiles, with M2 displaying larger mid-descent $\delta T$ and $\tau$ activity to steer onto the curved approach.

    \item \textit{Terminal sensitivity.} Perturbing the Phase~A$\to$B handoff state by $\pm0.05$ on each component (holding the optimal costate and $t_f$ fixed) raises the worst-case terminal residual to $\sim51$ for M1 but $\sim1.9\times10^3$ for M2. The sphere constraint and collinearity condition $p_{p_x}(t_f)\,p_z(t_f)=p_{p_z}(t_f)\,(p_x(t_f)-p_c)$ tightly couple position and costate at touchdown on M2.

    \item \textit{Feasibility.} Both scenarios satisfy $p_z\ge0$ along the trajectory, converge with $\lVert res\rVert_\infty\sim6.2\times10^{-8}$ (below the $10^{-5}$ target), and meet the terminal manifold and $H_B(t_f)=0$ residuals in Table~\ref{tab:val_p3}. Touchdown occurs at $(0,0)$ and $(0,5)$~m as intended.
\end{itemize}

\textit{Effect of manifold geometry.} Flat ground $M_1=\{p_z=0,\,v_z=0,\,\theta=0\}$ allows a nearly vertical descent over the shared footprint (Fig.~\ref{fig:p3_traj}, left). The platform manifold $M_2=\{h_1(x)=0,\,v_z=0,\,\theta=0\}$ replaces the altitude fix with the sphere $h_1=(p_x-p_c)^2+p_z^2-r^2=0$, forcing a bowed trajectory that intercepts the surface along its outward normal at the apex; transversality requires the position costate to align with that normal. The resulting arc demands sustained lateral thrust and attitude correction through Phase~B (Fig.~\ref{fig:p3_traj}, right), which---together with the higher $Q$ weights used for M2---explains the shorter but energetically costlier landing relative to M1.

\begin{figure}[H]
    \centering
    \includegraphics[width=\linewidth]{figures/p3_manifold_comparison.png}
    \caption{M1 vs M2 comparison: total landing time $t_f$, Phase~B duration, Phase~B control energy $\int u^\top R u\,dt$, Phase~A thrust switches, shooting feasibility $\lVert res\rVert_\infty$, and terminal sensitivity to $\pm0.05$ handoff perturbations.}
    \label{fig:p3_comparison}
\end{figure}

\subsubsection*{Validation}

All Part~III necessary conditions are verified to numerical tolerance: the terminal manifold and transversality residuals, the free-final-time condition $H_B(t_f)=0$, the Hamiltonian consistency $H_A\equiv0$ (minimum time) and $H_B\equiv0$ (free final time, autonomous), and the state constraint $p_z\ge0$ along the full trajectory. Table~\ref{tab:val_p3} reports the computed residuals.

\begin{table}[H]
\centering
\footnotesize
\caption{Part III Mission: PMP / TPBVP Validation Residuals}
\label{tab:val_p3}
\begin{tabular}{@{}lcc@{}}
\toprule
Metric & M1 (flat) & M2 (platform) \\ \midrule
Shooting $\lVert res\rVert_\infty$ (target $10^{-5}$) & $6.2\times10^{-8}$ & $6.2\times10^{-8}$ \\
Terminal manifold $\lVert res\rVert_\infty$ & $7.0\times10^{-21}$ & $1.1\times10^{-14}$ \\
Free-time condition $H_B(t_f)$ & $-7.0\times10^{-21}$ & $0.0$ \\
$H_A$ drift (min-time, $\equiv0$) & $2.9\times10^{-9}$ & $2.9\times10^{-9}$ \\
$H_B$ drift (free-time, $\equiv0$) & $9.7\times10^{-12}$ & $9.6\times10^{-11}$ \\
Min.\ altitude $p_z$ [m] ($\ge0$) & $0.000$ & $0.000$ \\ \midrule
Ascent time $t_1$ [s] & $2.158$ & $2.158$ \\
Landing time $t_f$ [s] & $8.097$ & $7.315$ \\
Phase~B energy $\int u^\top R u\,dt$ & $215.7$ & $526.2$ \\
Touchdown $(p_x,p_z)$ [m] & $(0,0)$ & $(0,5)$ \\
\bottomrule
\end{tabular}
\end{table}

\begin{figure}[H]
    \centering
    \includegraphics[width=\linewidth]{figures/p3_hamiltonian.png}
    \caption{Hamiltonian consistency: $H_A\equiv0$ on Phase~A (minimum time) and $H_B\equiv0$ on Phase~B (free final time), both to numerical tolerance.}
    \label{fig:p3_hamiltonian}
\end{figure}

%%%%%%%%%%%%%%%%%%%%%%%%%%%%%%%%%%%%%%%%%%%%%%%%%%%%%%%%%%%%%%%%%%%%%%%%%%%%%%%%
\newpage
\section{PART IV - The Nonlinear Optimal Control Near Hover}

\subsection{Phase Formulation}

Consider the fixed-time, free-endpoint nonlinear optimal control problem:

\[\min_{u} \int_{0}^{t_f} \left( x^\top Q x + u^\top R u \right) dt\]

Recall: 
State Equations (which is also $\dot{x} = \frac{\partial H}{\partial p}$):

\[
\dot{x} = \begin{bmatrix} 
v_x \\ 
v_z \\ 
-\frac{T}{m}\sin\theta \\
\frac{T}{m}\cos\theta - g \\ 
\omega \\ 
\frac{\tau}{I} \\ 
-\alpha T 
\end{bmatrix}
\]

System is anchored by the known initial launch state at $t=0$: 
\[x(0) = x_{0}\]
    
\subsection{Analytical Derivations}

\begin{itemize}
    \item Nonlinear Hamiltonian
    
    Adjoin the nonlinear dynamics to the running cost using the costate vector $p = [p_1, p_2, p_3, p_4, p_5, p_6, p_7]^\top$:

    \[
    \begin{aligned}
        H(x, u, p) ={}& x^\top Q x + u^\top R u + p_1 v_x + p_2 v_z - p_3 \tfrac{T}{m}\sin\theta \\
        &+ p_4 \big( \tfrac{T}{m}\cos\theta - g \big) + p_5 \omega + p_6 \tfrac{\tau}{I} - p_7 \alpha T
    \end{aligned}
    \]

    \item Adjoint equations (costate equations)
    \[\dot{p} = -\frac{\partial H}{\partial x}\]
    
    \[
    \begin{aligned}
    \dot{p}_1 &= -2Q_{11} p_x \\
    \dot{p}_2 &= -2Q_{22} p_z \\
    \dot{p}_3 &= -2Q_{33} v_x - p_1 \\
    \dot{p}_4 &= -2Q_{44} v_z - p_2 \\
    \dot{p}_5 &= -2Q_{55} \theta + p_3 \frac{T}{m}\cos\theta + p_4 \frac{T}{m}\sin\theta \\
    \dot{p}_6 &= -2Q_{66} \omega - p_5 \\
    \dot{p}_7 &= -2Q_{77} m - p_3 \frac{T}{m^2}\sin\theta + p_4 \frac{T}{m^2}\cos\theta
    \end{aligned}
    \]
    
    \item Stationarity conditions ($\frac{\partial H}{\partial u} = 0$)

    Consistent with the implementation and with Parts~I--II, the running cost penalizes the \emph{deviation} control $u=[\delta T,\ \tau]^\top$, where $\delta T = T - m_0 g$ is the thrust increment about the hover trim and the thrust entering the equations of motion is $T = m_0 g + \delta T$ (so $\partial T/\partial \delta T = 1$). Differentiating $H$ directly with respect to the deviation control yields the optimal increment, with no trim offset subtracted afterwards:

    \begin{itemize}
        \item For the thrust increment ($\delta T$):
        \[\frac{\partial H}{\partial \delta T} = 2R_{11}\,\delta T - p_3 \frac{\sin\theta}{m} + p_4 \frac{\cos\theta}{m} - p_7 \alpha = 0\]
        \[\delta T^* = \text{sat}_{\,\delta T_{\min}}^{\,\delta T_{\max}} \left[ \frac{1}{2 R_{11}} \left( p_3 \frac{\sin\theta}{m} - p_4 \frac{\cos\theta}{m} + p_7 \alpha \right) \right],\]
        where the saturation limits are the actuator bounds expressed in deviation form, $\delta T \in [T_{\min}-m_0 g,\ T_{\max}-m_0 g]$.

        \item For Torque ($\tau$):
        \[\frac{\partial H}{\partial \tau} = 2R_{22}\tau + \frac{p_6}{I} = 0\]
        \[\tau^* = \text{sat}_{\tau_{\min}}^{\tau_{\max}} \left[ -\frac{p_6}{2 I R_{22}} \right]\]
    \end{itemize}

    \item Transversality (terminal condition)
    
    Because the terminal state is subject to optimization, the terminal costate at $t=t_f$ is governed by the transversality condition based on the terminal penalty $\phi(x(t_f))$ (or specific manifold constraints):
    \[p(t_f) = \frac{\partial \phi}{\partial x} \bigg|_{t=t_f}\]

    \item Complete TPBVP
    
    The complete TPBVP is formed by substituting the optimal, saturated control laws ($\delta T^*$, $\tau^*$) back into the 7 state equations and 7 adjoint equations. This creates a deeply coupled system of 14 nonlinear, first-order ordinary differential equations. This system must be solved simultaneously across the domain $t \in [0, t_f]$, subject to the split boundary conditions $x(0)$ and $p(t_f)$.

    \begin{equation} \label{eq:full_tpbvp}
    \small
    \begin{aligned}
        % Forward State Dynamics
        \dot{p}_x &= v_x \\
        \dot{p}_z &= v_z \\
        \dot{v}_x &= -\frac{T}{m}\sin\theta \\
        \dot{v}_z &= \frac{T}{m}\cos\theta - g \\
        \dot{\theta} &= \omega \\
        \dot{\omega} &= \frac{\tau}{I} \\
        \dot{m} &= -\alpha T \\[1.5ex] 
        % Backward Costate Dynamics
        \dot{p}_1 &= -2Q_{11} p_x \\
        \dot{p}_2 &= -2Q_{22} p_z \\
        \dot{p}_3 &= -2Q_{33} v_x - p_1 \\
        \dot{p}_4 &= -2Q_{44} v_z - p_2 \\
        \dot{p}_5 &= -2Q_{55} \theta + \frac{T}{m} \big(p_3 \cos\theta + p_4 \sin\theta \big) \\
        \dot{p}_6 &= -2Q_{66} \omega - p_5 \\
        \dot{p}_7 &= -2Q_{77} m + \frac{T}{m^2} \big(p_4 \cos\theta - p_3 \sin\theta \big)
    \end{aligned}
    \end{equation}
    
\end{itemize}

\subsection{Numerical Implementation}

The Part~IV problem is solved in hover-deviation coordinates ($x$ is the deviation from the hover trim $T^\star = m_0 g$, $\theta^\star=0$; control $u=[\delta T,\ \tau]^\top$), so that the linearized solutions of the previous parts can serve directly as the nonlinear warm start. The implementation (\texttt{nonlinear.py}) builds three comparison scenarios on a common initial condition $x_0$ and a common short horizon:

\begin{itemize}
    \item \textbf{Restriction to the near-hover regime.} Because true hover is \emph{not} an equilibrium of a variable-mass vehicle (holding $T=mg$ continuously depletes mass at $\dot m = -\alpha m g$), the horizon is kept short, $t_f = 4\,\mathrm{s}$ with $201$ uniform nodes. This keeps the trajectory near hover and the mass safely above $m_{dry}$, which is the validity regime prescribed for this part.
    \item \textbf{Scenario 1 — LQR on the linear plant.} The backward Riccati sweep of Part~I produces $P(t)$ and the feedback $u=-R^{-1}B^\top P(t)\,x$, integrated forward on $\dot x = Ax+Bu$.
    \item \textbf{Scenario 2 — LQR on the nonlinear plant.} The \emph{same} time-varying gain is applied to the full nonlinear EOM, integrated with a fixed-step RK4 scheme, with control saturation enforced by clipping. This isolates the modeling error incurred by using the linear law on the true plant.
    \item \textbf{Scenario 3 — Nonlinear PMP / TPBVP.} The coupled $14$-state state/costate system of Eq.~\eqref{eq:full_tpbvp} is solved by collocation (\texttt{scipy.integrate.solve\_bvp}). The optimal control is substituted in closed form from the stationarity conditions and saturated to the actuator bounds. The split boundary conditions are the initial-state match $x(0)=x_0$ and the free-endpoint transversality $p(t_f) = 2 Q_f\, x(t_f)$.
    \item \textbf{Warm start and costate scaling.} The collocation guess uses the Scenario~1 state history and the costate guess $p(t) = 2P(t)x(t)$; the factor of two maps the $\tfrac12$-scaled Riccati costate to the no-$\tfrac12$ running cost $x^\top Qx + u^\top Ru$ used here. Convergence is declared at a collocation tolerance of $10^{-6}$.
\end{itemize}

\subsection{Analysis}

Three scenarios share $x_0=[1,\,1.5,\,0.4,\,-0.25,\,0.12,\,0,\,0.05]^\top$ and $t_f=4$~s: (1)~LQR on the LTI plant, (2)~the same LQR law on the nonlinear plant, and (3)~nonlinear PMP/collocation on the nonlinear plant with an LQR warm start.

\begin{figure}[H]
    \centering
    \includegraphics[width=\linewidth]{figures/p4_state_trajectories.png}
    \caption{State trajectories for the three scenarios.}
    \label{fig:p4_traj_comp}
\end{figure}

\begin{figure}[H]
    \centering
    \includegraphics[width=0.5\linewidth]{figures/p4_mission_plane_quiver.png}
    \caption{Mission plane $p_x$ vs.\ $p_z$ with attitude vectors.}
    \label{fig:p4_traj_comp2}
\end{figure}

\begin{figure}[H]
    \centering
    \includegraphics[width=\linewidth]{figures/p4_control_histories.png}
    \caption{Thrust and torque histories.}
    \label{fig:p4_control_hist}
\end{figure}

\begin{figure}[H]
    \centering
    \includegraphics[width=\linewidth]{figures/p4_energy_comparison.png}
    \caption{Control effort $J_u=\int u^\top R u\,dt$ and total cost $J$.}
    \label{fig:p4_energy}
\end{figure}

\begin{figure}[H]
    \centering
    \includegraphics[width=\linewidth]{figures/p4_terminal_behavior.png}
    \caption{Terminal state norm and components at $t_f$.}
    \label{fig:p4_term_behav}
\end{figure}

\begin{figure}[H]
    \centering
    \includegraphics[width=\linewidth]{figures/p4_convergence.png}
    \caption{Collocation convergence and tolerance sensitivity.}
    \label{fig:p4_conv_behav}
\end{figure}

\begin{figure}[H]
    \centering
    \includegraphics[width=\linewidth]{figures/p4_theta_sensitivity.png}
    \caption{Control effort and PMP savings vs.\ scaled initial pitch $\theta_0$.}
    \label{fig:p4_solver_sens}
\end{figure}

\begin{figure}[H]
    \centering
    \includegraphics[width=\linewidth]{figures/p4_terminal_residual.png}
    \caption{Per-component TPBVP boundary and optimality residuals.}
    \label{fig:p4_term_resid}
\end{figure}

Figs.~\ref{fig:p4_traj_comp}--\ref{fig:p4_traj_comp2} show scenarios~(1) and (2) agree early but diverge as $\theta$ and $m$ evolve; the nonlinear PMP solution trims peak $p_z$ and $\theta$ by exploiting $\sin\theta$ and $T/m$ coupling. The linear law over-commands on the true plant (Fig.~\ref{fig:p4_control_hist}), raising $J_u$ from $13.2$ to $42.9$ while PMP recovers $39.4$ ($\approx8\%$ savings) and lowers total $J$ by $41\%$ vs.\ scenario~(2) (Fig.~\ref{fig:p4_energy}). Terminal norms tighten under PMP (Fig.~\ref{fig:p4_term_behav}). Collocation converges in few refinements from the LQR seed (Fig.~\ref{fig:p4_conv_behav}); the PMP advantage grows with initial pitch (Fig.~\ref{fig:p4_solver_sens}). Boundary residuals meet the $10^{-6}$ target (Fig.~\ref{fig:p4_term_resid}). The linear approximation remains accurate for small attitude/mass deviations over $4$~s; it breaks down once thrust--gravity resolution and mass depletion accumulate.

\subsection{Validation}

To verify the numerical solution to the Two-Point Boundary Value Problem (TPBVP) and prove strict adherence to Pontryagin's Maximum Principle (PMP), the solution was validated against the necessary optimality conditions. 

\subsubsection*{Hamiltonian Consistency}
Part~IV is a \emph{fixed}-final-time, free-endpoint problem, and the hover-trim dynamics are autonomous. The Hamiltonian must therefore remain \emph{constant} along the optimal trajectory; unlike the free-final-time mission of Part~III, this constant is \emph{not} required to vanish. As shown in Figure \ref{fig:hamiltonian_phaseX2}, the Hamiltonian of the nonlinear PMP solution is constant to a drift of $\mathcal{O}(10^{-9})$ about $H(t_f) \approx 34.7$, confirming Hamiltonian consistency. By contrast, applying the linear LQR law to the nonlinear plant produces a Hamiltonian that drifts by $\mathcal{O}(10^{2})$ over the same horizon---a direct, quantitative certificate that the linear law is \emph{not} an extremal of the nonlinear problem.

\begin{figure}[H]
    \centering
    \includegraphics[width=\linewidth]{figures/p4_hamiltonian.png}
    \caption{Hamiltonian evolution $H(t)$ for the three scenarios. The nonlinear PMP solution (bottom) holds $H$ constant to $\mathcal{O}(10^{-9})$ as required for an autonomous fixed-time extremal, whereas the LQR law on the nonlinear plant drifts by $\mathcal{O}(10^{2})$.}
    \label{fig:hamiltonian_phaseX2}
\end{figure}

\subsubsection*{Terminal Conditions, Adjoint Boundaries, and Shooting Residuals}
Because the terminal state is free, the boundary conditions verified are the initial-state match $x(0)=x_0$ and the free-endpoint transversality $p(t_f) = 2 Q_f\, x(t_f)$; there are no fixed terminal-state targets in this part. Table \ref{tab:validation_phaseX} reports the per-component boundary residuals together with the interior collocation defect and the Hamiltonian drift.

The global $\mathcal{L}_\infty$ norm of the boundary-condition residual was $\mathbf{1.8 \times 10^{-15}}$ and the maximum interior collocation residual was $9.7 \times 10^{-7}$ on a converged mesh of $219$ nodes---both within the requested solver tolerance, confirming that the computed solution satisfies the necessary optimality conditions to numerical precision.

\begin{table}[H]
\centering
\caption{Part~IV validation: boundary and optimality residuals of the nonlinear TPBVP (fixed-time, free endpoint).}
\label{tab:validation_phaseX}
\begin{small} 
\setlength{\tabcolsep}{3.5pt} 
\begin{tabular}{@{}llc@{}}
\toprule
Type & Boundary / optimality condition & Residual \\ \midrule
Initial      & $x(0) = x_0$ \quad($\mathcal{L}_\infty$)                  & $3.9 \times 10^{-26}$ \\ \midrule
Transv.      & $p_{p_z}(t_f) = 2 Q_{f,22}\, p_z(t_f)$                     & $1.8 \times 10^{-15}$ \\
Transv.      & $p_{\theta}(t_f) = 2 Q_{f,55}\, \theta(t_f)$               & $8.9 \times 10^{-16}$ \\
Transv.      & $p_{\omega}(t_f) = 2 Q_{f,66}\, \omega(t_f)$               & $5.6 \times 10^{-17}$ \\
Transv.      & $p(t_f) = 2 Q_f\, x(t_f)$ \quad($\mathcal{L}_\infty$)      & $1.8 \times 10^{-15}$ \\ \midrule
Hamiltonian  & $H(t)$ drift (autonomous, fixed $t_f$)                     & $1.4 \times 10^{-9}$  \\
Collocation  & interior RMS defect (219 nodes)                           & $9.7 \times 10^{-7}$  \\ \bottomrule
\end{tabular}
\end{small}
\end{table}
 
\addtolength{\textheight}{-5cm}  

%%%%%%%%%%%%%%%%%%%%%%%%%%%%%%%%%%%%%%%%%%%%%%%%%%%%%%%%%%%%%%%
\section*{ACKNOWLEDGMENT}

The author thanks Professor Shahriar Talebi for guidance and instruction throughout the MAE 270C Optimal Control course. Their insights into optimal control theory formed the foundation of this project.

\begin{thebibliography}{99}

\bibitem{liberzon2012}
Liberzon, D. (2012). \emph{Calculus of variations and optimal control theory: A concise introduction}. Princeton University Press.

\bibitem{bryson2018}
Bryson, A. E., \& Ho, Y.-C. (2018). \emph{Applied optimal control}. Routledge. https://doi.org/10.1201/9781315137667

\bibitem{kirk2004}
Kirk, D. E. (2004). \emph{Optimal control theory: An introduction}. Courier Corporation.

\bibitem{acikmese2007}
Acikmese, B., \& Ploen, S. R. (2007). Convex programming approach to powered descent guidance for Mars landing. \emph{Journal of Guidance, Control, and Dynamics}, \emph{30}(5), 1353--1366. https://doi.org/10.2514/1.27553

\end{thebibliography}

\end{document}